%! Justifying a Claim About the Difference of Two Means Based on a CI — Unit 7, Lesson 7.7 — Lesson Plan
\documentclass[10pt]{article}
\usepackage{apstats-article}
\usepackage{apstats-boxes}
\usepackage{graphicx}
\graphicspath{{images/}}
\newcommand{\TallMath}[1]{$\displaystyle #1\rule[-1.4em]{0pt}{3.2em}$}
\newcommand{\UnitNumberName}{Unit 7: Inference for Quantitative Data: Means \quad}
\newcommand{\LessonNumberName}{Lesson 7.7: Justifying a Claim About the Difference of Two Means Based on a Confidence Interval}

\begin{document}
\begin{center}
    {\Large\bfseries \CourseName: \SchoolYear \\
    {\small\bfseries \UnitNumberName \LessonNumberName}}
\end{center}

\begin{tcolorbox}[colback=sky, colframe=navy, boxrule=0.9pt, arc=2mm,
    left=2mm, right=2mm, top=1.5mm, bottom=1.5mm]
    \textbf{Primary Objective:} Students will interpret a two-sample $t$-interval for
    $\mu_1 - \mu_2$ in context, use the interval to justify a claim (particularly whether
    it contains or excludes 0), and describe how sample size affects the width of the interval.
    (Big Idea: UNC)
\end{tcolorbox}

\renewcommand{\arraystretch}{1.6}
\begin{skillbox}[Priority Ideas \& Skills]{goldbox}
\begin{minipage}[t]{0.48\linewidth}
    \textbf{AP Skill 4B --- Interpretation}
    \begin{itemize}
        \item Interpret a confidence interval for $\mu_1 - \mu_2$ in context, referencing
              both samples and their populations (Skill 4.B; UNC-4.Z).
    \end{itemize}
    \textbf{AP Skill 4D --- Justification}
    \begin{itemize}
        \item Use a CI for $\mu_1 - \mu_2$ to justify or refute a claim; examine whether
              0 is inside or outside the interval (Skill 4.D; UNC-4.AA).
    \end{itemize}
    \textbf{AP Skill 4A --- Sample Size \& Width}
    \begin{itemize}
        \item Explain how increasing sample sizes decreases interval width when all other
              things remain the same (Skill 4.A; UNC-4.AB).
    \end{itemize}
\end{minipage}
\hfill
\begin{minipage}[t]{0.48\linewidth}
    \textbf{Key Understandings}
    \begin{itemize}
        \item A $C\%$ CI for $\mu_1 - \mu_2$ means: in repeated sampling with the same
              sizes, approximately $C\%$ of such intervals will capture the true difference.
        \item Interpretation must name both populations (not just samples) and give
              numbers in context with a direction (e.g., Group 1 $-$ Group 2).
        \item \textbf{If 0 is not in the interval:} plausible evidence of a real difference;
              the claim of ``no difference'' is not supported.
        \item \textbf{If 0 is in the interval:} a difference of 0 is plausible; insufficient
              evidence to claim the means differ.
        \item Larger $n$ $\Rightarrow$ smaller $SE$ $\Rightarrow$ narrower interval.
    \end{itemize}
\end{minipage}
\end{skillbox}

\begin{skillbox}[{Vocabulary, Concepts, \& Theorems}]{greenbox}
\begin{minipage}[t]{0.48\linewidth}
    \begin{itemize}
        \item \textbf{Two-sample $t$-interval for $\mu_1 - \mu_2$} --- interval estimate
              for the difference of two population means using $t^*$ with Welch df
        \item \textbf{Interpretation template} --- ``Based on these samples, we are $C\%$
              confident the true difference in population means ($\mu_1 - \mu_2$) is
              between [lower] and [upper] [units], in context.''
        \item \textbf{Plausible values} --- values inside the CI are consistent with the data
    \end{itemize}
\end{minipage}
\hfill
\begin{minipage}[t]{0.48\linewidth}
    \begin{itemize}
        \item \textbf{Zero in the interval} --- indicates a difference of 0 is plausible;
              cannot conclude the population means differ
        \item \textbf{Zero outside the interval} --- 0 is not a plausible difference;
              provides evidence the population means differ
        \item \textbf{Width of a CI} --- $2 \times t^* \cdot SE$; decreases as $n$ increases
              because $SE = \sqrt{s_1^2/n_1 + s_2^2/n_2}$ shrinks
        \item \textbf{Standard error (two-sample)} ---
              $SE = \sqrt{\dfrac{s_1^2}{n_1} + \dfrac{s_2^2}{n_2}}$
    \end{itemize}
\end{minipage}
\end{skillbox}

\begin{skillbox}[Activate Prior Knowledge \& Spiral Review]{sky}
\begin{minipage}[t]{0.55\linewidth}\vspace{0pt}
    \textbf{Spiral Review (warm-up)}
    \begin{itemize}
        \item Constructing a two-sample $t$-interval: formula, conditions, and Welch df
              --- from Lesson 7.6
        \item Interpreting a one-sample CI for $\mu$ in context --- Lesson 7.3
        \item Using a CI to justify a claim about a single proportion or mean --- Units 6--7
        \item Effect of sample size and confidence level on CI width --- Lessons 7.3, 7.6
    \end{itemize}
\end{minipage}
\hfill
\begin{minipage}[t]{0.38\linewidth}\vspace{0pt}\centering
    % Authored warm-up compiles to target/ — no source PDF to embed here.
\end{minipage}
\end{skillbox}

\begin{skillbox}[Hook]{sky}
\begin{minipage}[t]{0.55\linewidth}
    \textbf{Which City Has Faster Emergency Response? (3--4 min)}\\
    A journalist collects response-time data from two cities and produces this calculator
    output:
    \begin{center}
    \texttt{Two-Sample t Interval for $\mu_1 - \mu_2$}\\
    \texttt{95\% CI: ($-$2.37, 0.37) minutes}
    \end{center}
    Ask: ``Does this interval prove one city is faster?''\\[4pt]
    Key discussion points:
    \begin{itemize}
        \item The interval \emph{includes 0} --- a difference of zero is plausible.
        \item The data do not provide convincing evidence that the mean response times differ.
        \item This is the 2009 AP FRQ 4 illustrative example (UNC-4.Z.2).
    \end{itemize}
\end{minipage}
\hfill
\begin{minipage}[t]{0.40\linewidth}
    \textbf{Bridge from Lesson 7.6}\\
    In Lesson 7.6 we \emph{built} the interval:
    \[(\bar x_1 - \bar x_2) \pm t^* \cdot \sqrt{\tfrac{s_1^2}{n_1} + \tfrac{s_2^2}{n_2}}\]
    Today we \emph{read} and \emph{use} the interval to answer a question:\\[4pt]
    \begin{center}
    \textbf{Is 0 inside or outside?}\\
    Inside $\to$ no convincing evidence of a difference.\\
    Outside $\to$ evidence of a real difference.
    \end{center}
\end{minipage}
\end{skillbox}

\begin{skillbox}[Lesson: Interpreting a Two-Sample CI]{sky}
\begin{minipage}[t]{0.48\linewidth}
    \textbf{The Interpretation Template (UNC-4.Z.2)}\\
    An interpretation must include:
    \begin{enumerate}
        \item The confidence level ($C\%$)
        \item Both populations (not the samples)
        \item The direction of subtraction ($\mu_1 - \mu_2$)
        \item The lower and upper bounds with units
    \end{enumerate}
    \textit{Model sentence:}\\
    ``Based on these samples, we are 95\% confident that the difference in population mean
    response times (northern $-$ southern) is between $-2.37$ minutes and $0.37$ minutes.''
\end{minipage}
\hfill
\begin{minipage}[t]{0.48\linewidth}
    \textbf{Connecting to the Long-Run Frequency (UNC-4.Z.1)}\\
    \begin{itemize}
        \item The 95\% refers to the \emph{method}, not to this specific interval.
        \item In repeated random sampling with the same sample sizes, approximately 95\% of
              all constructed intervals will contain $\mu_1 - \mu_2$.
        \item We do \emph{not} say ``there is a 95\% probability that $\mu_1 - \mu_2$
              falls in this interval'' (the parameter is fixed; the interval is random).
    \end{itemize}
\end{minipage}
\end{skillbox}

\begin{skillbox}[Lesson (cont.): Justifying a Claim \& Width]{sky}
\begin{minipage}[t]{0.48\linewidth}
    \textbf{Using the CI to Justify a Claim (UNC-4.AA)}
    \renewcommand{\arraystretch}{1.5}
    \begin{center}
    \begin{tabular}{@{}p{0.40\linewidth} p{0.52\linewidth}@{}}
        \textbf{Situation} & \textbf{Conclusion} \\ \hline
        0 is \emph{inside} the interval & A difference of 0 is plausible; insufficient
          evidence to claim $\mu_1 \ne \mu_2$. \\
        0 is \emph{outside} the interval (entire interval $>0$) & Plausible evidence
          $\mu_1 > \mu_2$; supports a ``positive difference'' claim. \\
        0 is \emph{outside} the interval (entire interval $<0$) & Plausible evidence
          $\mu_1 < \mu_2$; supports a ``negative difference'' claim. \\
    \end{tabular}
    \end{center}
\end{minipage}
\hfill
\begin{minipage}[t]{0.48\linewidth}
    \textbf{Effect of Sample Size on Width (UNC-4.AB)}
    \begin{itemize}
        \item $SE = \sqrt{\dfrac{s_1^2}{n_1} + \dfrac{s_2^2}{n_2}}$ decreases as
              $n_1$ and/or $n_2$ increase.
        \item Smaller $SE$ $\Rightarrow$ narrower interval $\Rightarrow$ more precision.
        \item Doubling both sample sizes does \emph{not} halve the width (only
              $1/\sqrt{2} \approx 0.71\times$ narrower).
        \item Higher confidence level ($C\%$ $\uparrow$) $\Rightarrow$ larger $t^*$
              $\Rightarrow$ wider interval (all else equal).
    \end{itemize}
\end{minipage}
\end{skillbox}

\begin{skillbox}[Active Monitoring]{redbox}
\begin{minipage}[t]{0.48\linewidth}
    \textbf{Circulate and Check}
    \begin{itemize}
        \item Interpretation names \emph{populations}, not samples; includes both bounds
              and units; specifies direction of subtraction.
        \item Claim justification explicitly checks whether 0 is inside or outside.
        \item Width reasoning references $SE$ (and sample sizes), not just the formula.
        \item Students do not confuse ``95\% confident'' with ``95\% probability the
              parameter is in the interval.''
    \end{itemize}
\end{minipage}
\hfill
\begin{minipage}[t]{0.48\linewidth}
    \textbf{Cold-Call Prompts}
    \begin{itemize}
        \item ``Why do we need to say 'population means' and not just 'means'?''
        \item ``The interval is $(1.2, 4.8)$ minutes. What does that tell us about the
              claim that the two groups have the same mean?''
        \item ``If we doubled both sample sizes, what would happen to the interval?''
        \item ``Does 95\% confident mean there's a 95\% chance the true value is inside?
              Explain the correct meaning.''
    \end{itemize}
\end{minipage}
\end{skillbox}

\begin{skillbox}[Group Work \& Differentiation]{redbox}
\begin{multicols}{3}
    \textbf{Tier R --- Remediate}
    \begin{itemize}
        \item Provide a fill-in sentence frame for interpretation.
        \item Focus on locating 0 relative to the interval and choosing the correct
              conclusion word (``supports'' vs.\ ``does not support'').
        \item Activity: Tier R template task only.
    \end{itemize}
    \columnbreak
    \textbf{Tier A --- Approaching}
    \begin{itemize}
        \item Calculate a two-sample CI from summary statistics, write a full
              interpretation, and answer a claim question with justification.
        \item Explain in words why a larger sample narrows the interval.
    \end{itemize}
    \columnbreak
    \textbf{Tier E --- Extension}
    \begin{itemize}
        \item Compare two scenarios differing only in sample size; quantify the
              difference in width.
        \item Design a study: determine minimum $n$ needed to achieve a desired
              margin of error (conceptual reasoning).
    \end{itemize}
\end{multicols}
\end{skillbox}

\begin{skillbox}[Individual Work \& Assessment]{redbox}
\begin{minipage}[t]{0.48\linewidth}
    \textbf{Exit Ticket (5 min)}\\
    A 90\% CI for $\mu_{\text{urban}} - \mu_{\text{rural}}$ sleep hours is $(0.3, 1.1)$ hr.
    \begin{enumerate}
        \item Interpret the interval in context.
        \item Does the interval provide convincing evidence that urban residents sleep
              more on average? Explain.
    \end{enumerate}
\end{minipage}
\hfill
\begin{minipage}[t]{0.48\linewidth}
    \textbf{AP-Style Check}\\
    \textit{A 95\% CI for $\mu_A - \mu_B$ battery life is $(-3.2, 0.4)$ hours.
    Which conclusion is most appropriate?}
    \begin{enumerate}[label=(\Alph*)]
        \item Battery A lasts longer; the interval is mostly negative.
        \item There is convincing evidence that the means differ because the interval
              is not symmetric about 0.
        \item There is insufficient evidence of a difference because 0 is in the interval.
              \textcolor{keyred}{\textbf{$\leftarrow$ correct}}
        \item Battery B lasts longer because the lower bound is negative.
    \end{enumerate}
\end{minipage}
\end{skillbox}

\begin{skillbox}[Reinforcement \& Extension]{goldbox}
\begin{minipage}[t]{0.48\linewidth}
    \textbf{Homework Overview}
    \begin{itemize}
        \item Four problems: interpret given CIs for $\mu_1 - \mu_2$, justify claims, and
              analyze the effect of sample size on width.
        \item One problem revisits the 2009 AP FRQ 4 format (northern vs.\ southern
              response times).
    \end{itemize}
\end{minipage}
\hfill
\begin{minipage}[t]{0.48\linewidth}
    \textbf{Extension / Preview}
    \begin{itemize}
        \item \textit{Extension:} If a 95\% CI excludes 0, must a corresponding
              two-sided significance test at $\alpha = 0.05$ also reject $H_0$?
              Explain the connection.
        \item \textit{Preview:} Lesson 7.8 introduces significance tests for
              $\mu_1 - \mu_2$ --- the two-sample $t$-test.
    \end{itemize}
\end{minipage}
\end{skillbox}

\end{document}
